% Pista geo-26j2-q4b · Geometria
% Procedència: PAU Catalunya 2026 · Convocatòria de juny · Sèrie 5 · Exercici 4 — Opció B
\documentclass[11pt,a4paper]{article}
\usepackage{pista}
\pistainfo{Catalunya 2026}{Convocatòria de juny}{Sèrie 5}

\begin{document}

\section*{\textcolor{blauPista}{Exercici 4 --- Opció B}}

\noindent En una indústria es fan servir braços robòtics amb sensors per a mesurar amb precisió determinades peces metàl·liques. Les peces es col·loquen sobre la superfície determinada pel pla $\pi: (x, y, z) = (6, 1, 1) + \lambda(1, 1, 0) + \mu(0, 0, 1)$ i, per damunt seu, un braç robòtic desplaça un sensor des del punt $A = (1, 2, 3)$ fins al punt $B = (5, 3, 7)$ en línia recta.

\subsection*{\textit{a)} \normalfont Quina distància ha recorregut el sensor? \hfill\textnormal{[0,5 punts]}}

\pista{El desplaçament del sensor és, simplement, el mòdul del vector $\vec{AB}$.}

\subsection*{\textit{b)} \normalfont En acabar el moviment, a quina distància de la superfície es troba el sensor? \hfill\textnormal{[1 punt]}}

\pista{Et donen el pla $\pi$ en forma vectorial (paramètrica): abans de res, passa'l a la forma general $Ax + By + Cz + D = 0$ (per exemple, calculant el vector normal amb el producte vectorial dels dos vectors directors, $(1, 1, 0)$ i $(0, 0, 1)$, i imposant que el pla passi pel punt $(6, 1, 1)$). Un cop tinguis l'equació general, aplica la fórmula de la distància d'un punt --el punt $B$, que és on acaba el sensor-- a un pla.}

\subsection*{\textit{c)} \normalfont Si el sensor seguís movent-se en la mateixa línia recta i en el mateix sentit, en quin punt xocaria amb la superfície? \hfill\textnormal{[1 punt]}}

\pista{Escriu l'equació paramètrica de la recta que passa per $A$ i $B$ (amb vector director $\vec{AB}$). El ``xoc'' amb la superfície és, simplement, el punt d'intersecció d'aquesta recta amb el pla $\pi$: substitueix les coordenades paramètriques de la recta a l'equació general del pla (la que has trobat a l'apartat anterior) i resol per trobar el valor del paràmetre; amb aquest valor, calcula les coordenades concretes del punt.}

\end{document}
